\documentclass[10pt,a4paper,landscape]{article}
\usepackage[a4paper,landscape,top=1.5cm,bottom=1.8cm,left=1.6cm,right=1.6cm]{geometry}
\usepackage{booktabs}
\usepackage{array}
\usepackage{longtable}
\usepackage{tabularx}
\usepackage{xcolor}
\usepackage{titlesec}
\usepackage[hidelinks]{hyperref}
\usepackage{amsmath}

\renewcommand{\arraystretch}{1.15}

% Column types for the two tables
\newcolumntype{V}{>{\ttfamily\raggedright\arraybackslash}p{1.5cm}}       % variant / codon in typewriter
\newcolumntype{W}{>{\ttfamily\raggedright\arraybackslash}p{1.1cm}}       % narrow typewriter
\newcolumntype{T}{>{\ttfamily\raggedright\arraybackslash}p{3.0cm}}       % wider typewriter for target codon
\newcolumntype{N}{>{\raggedright\arraybackslash}X}                       % notes - stretch to fill
\newcolumntype{S}{>{\raggedright\arraybackslash}p{3.0cm}}                % source
\newcolumntype{R}{>{\raggedright\arraybackslash}p{2.2cm}}                % reachability
\newcolumntype{G}{>{\raggedright\arraybackslash}p{3.4cm}}                % signature concordance

\titleformat{\section}{\normalfont\large\bfseries}{\thesection}{0.5em}{}

\begin{document}

\begin{center}
{\Large\bfseries Supplementary Table S1\\[2pt]}
{\normalsize Per-variant provenance and codon reachability for the KRAS G12C and ALK composite resistance panels}
\end{center}

\vspace{4pt}

\noindent\textbf{Companion to:} Ronco, M. \emph{Quantifying the preclinical blind spots of mutagen-based TKI-resistance screens from chemistry-free saturation mutagenesis atlases.} 2026.

\vspace{10pt}

% ============================================================================
% S1a — KRAS G12C
% ============================================================================
\section*{S1a. KRAS G12C resistance panel (12 variants)}

\noindent Codon numbering follows ClinVar records anchored to NM\_033360.4 (KRAS-4A). The wild-type codon reported here is the codon in the \emph{G12C-mutant background} that is the shared starting point for every preclinical screen in the Awad 2021, Feng 2022, Koga 2021, Georgakopoulos-Soares 2025, and Tang 2022 (Boehringer) studies cited in the main text. Codon 12 is therefore TGT (not GGT); all other panel codons are wild-type at the KRAS reference. Reachability is classified by Hamming distance from the wild-type codon to any codon encoding the target amino acid: \emph{two-step SNV} if the minimum distance is $\geq 2$, meaning the variant is suppressed rather than forbidden (main text \S2.2); \emph{multi-class SNV} if a single-nucleotide route exists but more than one pyrimidine substitution class produces the amino acid; \emph{class-forced} if a single-nucleotide route exists and only one pyrimidine substitution class produces the amino acid.

\vspace{6pt}

\noindent\footnotesize
\begin{tabularx}{\textwidth}{V S W T R c G N}
\toprule
\textbf{Variant} & \textbf{Source panel} & \textbf{WT codon} & \textbf{Target codon} & \textbf{Reachability} & \textbf{Class} & \textbf{Signature concordance} & \textbf{Notes} \\
\midrule
C12F & Awad 2021 \emph{in vitro} DMS & TGT & TTT (pos 2 G$>$T) & class-forced & C$>$A & SBS4 (tobacco) & G12C-mutant background; single-nt via pos 2. \\
V8L  & Awad 2021 \emph{in vitro} DMS & GTA & CTA / TTA         & multi-class SNV & --- & multi-class SNV & Reachable in one nt via C$>$G (CTA) or C$>$A (TTA). \\
V8M  & Awad 2021 \emph{in vitro} DMS & GTA & ATG (requires 2 nt) & two-step SNV & --- & two-step, 1 route & Two substitutions required (pos 1 G$>$A and pos 3 A$>$G); 51- to 169-fold suppressed. \\
R68L & Awad 2021 \emph{in vitro} DMS & AGG & all Leu codons need $\geq 2$ nt & two-step SNV & --- & two-step, 2 routes & Top-ranked variant in Feng chemistry-free atlas under both drugs; 51- to 169-fold suppressed. \\
R68M & Awad 2021 \emph{in vitro} DMS & AGG & ATG (pos 2 G$>$T) & class-forced & C$>$A & SBS4 (tobacco) & Single-nt route via pos 2. \\
R68S & Awad 2021 clinical cohort     & AGG & AGT / AGC         & multi-class SNV & --- & multi-class SNV & C$>$A (AGT) or C$>$G (AGC) at pos 3. \\
H95D & Awad 2021 clinical cohort     & CAT & GAT (pos 1 C$>$G) & class-forced & C$>$G & SBS13 at TpC; else rare & H95D under adagrasib: chem-free rank 12 of 1,205, Mutazyme-simulated rank 90 (shift $+78$ at $\beta=1$; $+0$ to $+812$ across $\beta \in [0.25, 4]$). \\
H95Q & Awad 2021 clinical cohort     & CAT & CAA / CAG         & multi-class SNV & --- & multi-class SNV & T$>$A (CAA) or T$>$G (CAG) at pos 3. \\
H95R & Awad 2021 clinical cohort     & CAT & CGT (pos 2 A$>$G) & class-forced & T$>$C & SBS5 (clock-like) & Single-nt route via pos 2. \\
Y96C & Awad 2021 clinical cohort     & TAT & TGT (pos 2 A$>$G) & class-forced & T$>$C & SBS5 (clock-like) & Chemistry-free rank 2 under sotorasib, Mutazyme-simulated rank 2 (near-zero shift). \\
Y96D & Tanaka 2021 case report       & TAT & GAT (pos 1 T$>$G) & class-forced & T$>$G & low in common NSCLC sigs & First reported clinically in Tanaka 2021; chem-free rank 9 under sotorasib, Mutazyme-simulated rank 223 (shift $+214$ at $\beta=1$) --- class T$>$G is 4\%--11\% of ENU output and 3.6\% of Mutazyme II. \\
Y96H & Awad 2021 \emph{in vitro} DMS & TAT & CAT (pos 1 T$>$C) & class-forced & T$>$C & SBS5 (clock-like) & Single-nt route via pos 1. \\
\bottomrule
\end{tabularx}

\vspace{4pt}
\noindent\footnotesize
\textbf{Summary:} 2 two-step SNV (R68L, V8M); 3 multi-class SNV (V8L, R68S, H95Q); 7 class-forced (C12F, R68M, H95D, H95R, Y96C, Y96D, Y96H).

\vspace{18pt}
\normalsize

% ============================================================================
% S1b — ALK
% ============================================================================
\section*{S1b. ALK resistance panel (12 variants)}

\noindent Codon numbering follows the ALK NM\_004304.5 CDS cross-checked against the Oh 2026 Table S11 EML4-ALK v1 WT CDS. Reachability definitions as in S1a.

\vspace{6pt}

\noindent\footnotesize
\begin{tabularx}{\textwidth}{V S W T R c G N}
\toprule
\textbf{Variant} & \textbf{Source} & \textbf{WT codon} & \textbf{Target codon} & \textbf{Reachability} & \textbf{Class} & \textbf{Signature concordance} & \textbf{Notes} \\
\midrule
V1180L & Doi 2022 (alectinib)           & GTT & CTT (pos 1 G$>$C)        & class-forced    & C$>$G & SBS13 at TpC; else rare & Chem-free rank 8, Mutazyme-simulated rank 15 under alectinib (class C$>$G at 5\% in Mutazyme). \\
L1196M & Doi 2022 (alectinib)           & CTG & ATG (pos 1 C$>$A)        & class-forced    & C$>$A & SBS4 (tobacco)          & Gatekeeper. Rank 1 under both chem-free and Mutazyme-simulated rankings for alectinib. \\
I1171N & Doi 2022 (alectinib)           & ATC & AAC (pos 2 T$>$A)        & class-forced    & T$>$A & low in common NSCLC sigs & Compound-mutation background for the Takei 2025 neladalkib fixed-V1 screen. \\
I1171T & Doi 2022 (alectinib)           & ATC & ACC (pos 2 T$>$C)        & class-forced    & T$>$C & SBS5 (clock-like)        & Class T$>$C --- well-supplied by ENU and Mutazyme. \\
L1198F & Doi 2022 (alectinib)           & CTC & TTC (pos 1 C$>$T)        & class-forced    & C$>$T & SBS1 / SBS2 / SBS5 (common) & Class C$>$T --- well-supplied by Mutazyme II (47\% of substitutions). \\
G1202R & Doi 2022 (alectinib)           & GGA & CGA (C$>$G) / AGA (C$>$T) & multi-class SNV & ---   & multi-class SNV          & Two single-nt routes to Arg via different classes; not class-forced. \\
G1269A & Takei 2025 (next-gen ALK)      & GGA & GCA (pos 2 G$>$C)        & class-forced    & C$>$G & SBS13 at TpC; else rare  & Single-nt route via pos 2. \\
F1174C & Takei 2025 (next-gen ALK)      & TTC & TGC (pos 2 T$>$G)        & class-forced    & T$>$G & low in common NSCLC sigs & Pre-existing in a Lin 2024 ALKOVE-1 NVL-655 responder (cis with G1202R). \\
F1174L & Takei 2025 (next-gen ALK)      & TTC & CTC / TTA / TTG          & multi-class SNV & ---   & multi-class SNV          & Three single-nt routes to Leu (classes T$>$C, C$>$A, C$>$G); not class-forced. \\
L1196Q & Takei 2025 (next-gen ALK)      & CTG & CAG (pos 2 T$>$A)        & class-forced    & T$>$A & low in common NSCLC sigs & Compound-mutation partner reported with the gatekeeper residue. \\
C1156Y & Takei 2025 (next-gen ALK)      & TGC & TAC (pos 2 G$>$A)        & class-forced    & C$>$T & SBS1 / SBS2 / SBS5 (common) & Class C$>$T --- well-supplied by Mutazyme. \\
L1152R & Takei 2025 (next-gen ALK)      & CTG & CGG (pos 2 T$>$G)        & class-forced    & T$>$G & low in common NSCLC sigs & Largest predicted negative rank shift under Mutazyme $\times$ alectinib ($+558$ at $\beta=1$; $+115$ to $+817$ across $\beta \in [0.25, 4]$). \\
\bottomrule
\end{tabularx}

\vspace{4pt}
\noindent\footnotesize
\textbf{Summary:} 10 class-forced (V1180L, L1196M, I1171N, I1171T, L1198F, G1269A, F1174C, L1196Q, C1156Y, L1152R); 2 multi-class SNV (G1202R, F1174L). No two-step variants in the ALK panel.

\vspace{18pt}
\normalsize

% ============================================================================
% Caption / methodological notes
% ============================================================================
\section*{Methodological notes}

\footnotesize

\noindent\textbf{Panel construction.} The KRAS G12C panel is a \emph{composite} of patient-observed variants (Awad 2021 clinical cohort; Tanaka 2021 case report) and preclinical \emph{in vitro} DMS variants (Awad 2021 accompanying deep mutational scan). The ALK panel is drawn from Doi 2022 \emph{Targeted Oncology} (six variants recovered in the alectinib Mutazyme II Ba/F3 screen) and Takei 2025 \emph{Pharmaceutical Research} (six single-amino-acid variants from the next-generation ALK TKI screens). Panel membership is treated as fixed input to the analysis; the chemistry-free atlas rankings and class-forced-codon reachability classifications are derived.

\vspace{4pt}

\noindent\textbf{Doi 2022 ALK variants.} Doi 2022 reports resistance mutations in Figure 3a of the Mutazyme II $\times$ alectinib Ba/F3 screen without an explicit numeric ranking. The six Doi variants listed above are the single-amino-acid resistance mutations identified in Fig.~3a after excluding (i) silent passenger mutations (L1152 silent, L1165 silent, R1192 silent, all flagged in the Doi 2022 figure) and (ii) the I1171S sub-variant of the I1171N/T/S compound call (included in Doi but not in our panel because the three sub-variants cannot be distinguished at the position-level NGS readout; we retain I1171N and I1171T as the two independently reported sub-variants).

\vspace{4pt}

\noindent\textbf{Takei 2025 ALK variants.} The six Takei 2025 single-amino-acid variants listed above are drawn from the Takei \emph{Pharmaceutical Research} manuscript's resistance panel across its alectinib, lorlatinib, zotizalkib, and neladalkib screens. The Takei 2025 fixed-V1 compound-mutation screen reports additional second-step partners (L1256F, E1129V, E1129K, C1156Y, T1151R selected at 500 nM neladalkib on an I1171N background); these are discussed in the main text as compound-mutation predictions and are not included in the single-amino-acid resistance panel for Figure~3.

\vspace{4pt}

\noindent\textbf{Reachability classification.} Reachability is determined from the Hamming distance between the wild-type codon (from the WT CDS cross-checked against ClinVar nucleotide-level records; see data/reference/codons.yaml and codons\_kras.yaml) and the set of codons encoding the target amino acid. The classification is strictly a property of the genetic code and the wild-type codon; it does not depend on which mutagen chemistry a screen uses. The forced class, when a variant is class-forced, is the pyrimidine substitution class of the required single-nucleotide substitution: C$>$A, C$>$G, C$>$T, T$>$A, T$>$C, or T$>$G (all other substitutions are their complements and share a class with their pyrimidine representative).

\vspace{4pt}

\noindent\textbf{Chemistry-free rankings.} For KRAS G12C the chemistry-free ranking is the Z-score ranking of the 1,205 class-forced variants in the Feng 2022 Twist Biosciences synthetic saturation library under sotorasib or adagrasib. For ALK the chemistry-free ranking is $\sigma_{\mathrm{drug}}(v) = \text{Resistance}(v)$ (the Oh 2026 atlas Resistance score, used without further transformation --- see main-text Methods for the rationale, which turns on Oh's Resistance and Fitness scores being orthogonal axes rather than components to be subtracted) over the 1,259 class-forced variants in the Oh 2026 prime-editing saturation library under alectinib, lorlatinib, or zotizalkib. Per-variant scores and ranks are available in the derived data files \texttt{feng\_awad2021\_variant\_ranks.csv} and \texttt{oh\_decomposition\_all\_drugs.csv} in the paper's accompanying code repository.

\vspace{4pt}

\noindent\textbf{Signature concordance column.} Each class-forced variant's forced pyrimidine substitution class is annotated with the published COSMIC v3.3 single-base substitution (SBS) signatures (Alexandrov et al.\ \emph{Nature} 2020;578:94--101, DOI 10.1038/s41586-020-1943-3) that emphasise that class, as a qualitative tumour-context tag. The mapping is: C$>$A $\rightarrow$ SBS4 (tobacco smoking); C$>$T $\rightarrow$ SBS1 (CpG deamination), SBS2 (APOBEC at TpC), SBS5 (clock-like unknown); C$>$G $\rightarrow$ SBS13 (APOBEC at TpC); T$>$C $\rightarrow$ SBS5; T$>$A and T$>$G $\rightarrow$ low in all common NSCLC signatures. The two cohort contexts of this paper are not symmetric: KRAS G12C is a predominantly smoking-associated NSCLC population (SBS4 dominates), so C$>$A variants are maximally tumour-reachable and C$>$G/T$>$G variants are relatively rare somatically; ALK-rearranged LUAD is predominantly a never-smoker population (SBS1 and SBS5 dominate), so C$>$T is the maximally tumour-reachable class and SBS4-emphasised C$>$A variants (such as L1196M) are rarer somatically than their resistance phenotype would suggest. Variants flagged ``low in common NSCLC signatures'' are predicted to be acquired at low frequency somatically in both cohorts, so a mutagen-chemistry blind-spot in those classes has lower clinical priority than a blind-spot in a tumour-reachable class --- the \emph{tumour-reachable $\times$ mutagen-unreachable} intersection is the operationally sharp blind-spot. The 96-channel context-aware version of this analysis, with explicit per-signature per-trinucleotide weights and a per-cohort patient-level reweighting, is deferred to a companion paper that requires GENIE patient-level 96-channel catalogues not integrated here.

\end{document}
